\documentclass[11pt]{article}

\usepackage[T1]{fontenc}
\usepackage[utf8]{inputenc}
\usepackage{lmodern}
\usepackage[margin=1.02in]{geometry}
\usepackage{amsmath,amssymb,amsthm,mathtools}
\usepackage{booktabs}
\usepackage{array}
\usepackage{enumitem}
\usepackage{xcolor}
\usepackage{hyperref}
\usepackage{xurl}

\hypersetup{
  colorlinks=true,
  linkcolor=blue!55!black,
  citecolor=green!35!black,
  urlcolor=blue!55!black
}
\setlist{itemsep=2pt,topsep=4pt}
\setlength{\parindent}{0pt}
\setlength{\parskip}{5pt}
\setlength{\emergencystretch}{3em}

\newtheorem{theorem}{Theorem}[section]
\newtheorem{proposition}[theorem]{Proposition}
\newtheorem{lemma}[theorem]{Lemma}
\newtheorem{corollary}[theorem]{Corollary}
\newtheorem{definition}[theorem]{Definition}
\newtheorem{remark}[theorem]{Remark}
\newtheorem{warning}[theorem]{Warning}

\newcommand{\R}{\mathbb R}
\newcommand{\C}{\mathbb C}
\newcommand{\Q}{\mathbb Q}
\newcommand{\Z}{\mathbb Z}
\newcommand{\N}{\mathbb N}
\newcommand{\T}{\mathbb T}
\newcommand{\HGAS}{\mathrm{HGAS}}
\newcommand{\dd}{\,\mathrm d}
\newcommand{\Lean}{\textsc{Lean}}
\newcommand{\Mathlib}{\textsc{Mathlib}}
\newcommand{\RE}{\mathrm{RE}}
\newcommand{\GAS}{\mathrm{GAS}}
\newcommand{\leanname}[1]{\texttt{\detokenize{#1}}}
\newcommand{\leanmodule}[1]{\path{#1}}

\title{\textbf{Fixed-Resource Undecidability of Homogeneous Polynomial
Global Asymptotic Stability}\\
\large A detailed mathematical proof and guide to the complete Lean
formalization}
\author{Formalization project companion}
\date{24 July 2026}

\begin{document}
\maketitle

\begin{abstract}
For fixed natural numbers \(n,d\), let
\(\HGAS_{\Q}(n,d)\) be the set of finite sparse rational codes for
homogeneous polynomial vector fields \(F:\R^n\to\R^n\) of degree \(d\)
whose origin is globally asymptotically stable.  Invalid homogeneity codes
are declared negative instances.

This document proves that there are fixed constants \(n_\ast\geq1\) and
odd \(d_\ast\geq3\) for which
\(\HGAS_{\Q}(n_\ast,d_\ast)\) is
\(\Sigma^0_1\)-complete under computable many-one reductions.  In
particular, the fixed-pair decision problem is undecidable and its
complement is not computably enumerable.  More strongly, there is a fixed
threshold \(N_\ast\) such that
\[
  \HGAS_{\Q}(N,d)
  \quad\hbox{is \(\Sigma^0_1\)-complete whenever}\quad
  N\geq N_\ast,\quad d\geq3,\quad d\ \hbox{odd}.
\]
Thus cubic homogeneous GAS is already \(\Sigma^0_1\)-complete in every
fixed dimension \(N\geq N_\ast\).

The lower bound compiles universal halting into a fixed diagonal--quadratic--
linear integer system, a prime-radical torus detector, a polynomial
modular/Jacobi source, a compact stereographic accumulator, and finally a
rational homogeneous polynomial field.  The upper bound enumerates exact
rational Lyapunov certificates.  A solvability-preserving augmentation
\(x_j-y_j=0\) is essential: it forces every detector column to have a
nonzero linear character and thereby removes a genuine zero-ratio
resonance in an earlier three-moment argument.

Every theorem used in the final dependency chain is formalized in the
accompanying \Lean{} project.
The final theorems contain no \leanname{sorry}, \leanname{admit}, project
\leanname{axiom}, or \leanname{unsafe} declaration.  Kernel inspection of
the headline results reports only \leanname{propext},
\leanname{Classical.choice}, and \leanname{Quot.sound}.  The fixed constants
are named finite objects selected once and for all; because the particular
MRDP representative is obtained by classical choice, the development does
not claim that \(n_\ast,d_\ast,N_\ast\) have been evaluated to decimal
numerals.
\end{abstract}

\tableofcontents

\section{The theorem, exactly stated}
\label{sec:headline}

\subsection{Finite syntax and semantics}

A \emph{polynomial field code} in dimension \(n\) is an \(n\)-tuple of
finite lists of pairs
\[
    (\alpha,c),\qquad
    \alpha\in\N^n,\quad c\in\Q.
\]
Its realization is the rational polynomial vector field
\[
 F_i(x)=\sum_{(\alpha,c)\text{ in coordinate }i}c\,x^\alpha.
\]
Duplicate monomials and zero coefficients are allowed in the raw syntax;
decoding performs the usual finite polynomial sum.  There is a primitive
recursive Boolean test saying that every decoded coordinate is homogeneous
of degree \(d\).

The formal predicate is
\[
\begin{split}
 \HGAS_{\Q}(n,d)(F)
 \quad\Longleftrightarrow\quad&
 F\text{ passes the degree-\(d\) homogeneity test},\\
 &0\text{ is globally asymptotically stable for }\dot x=F(x).
\end{split}
\tag{1.1}
\]
This is \leanname{HGASQ} in
\leanmodule{HGAS/Computability/HGASLanguage.lean}.  Global asymptotic
stability is not defined through a possibly nonexistent global flow.  It
contains four separate assertions:

\begin{enumerate}
\item the origin is an equilibrium;
\item every initial point admits a forward solution on \([0,\infty)\);
\item the origin is Lyapunov stable; and
\item every forward solution converges to the origin.
\end{enumerate}

These definitions occur in
\leanmodule{HGAS/Dynamics/Stability.lean}.  This explicit forward-completeness
clause matters because polynomial ODEs may otherwise blow up in finite
time.

\begin{definition}
A predicate \(P\) on a primitively codable type is
\(\Sigma^0_1\)-complete if \(P\) is computably enumerable and every
computably enumerable predicate on \(\N\) computably many-one reduces to
\(P\).
\end{definition}

The project uses \Mathlib{}'s partial-recursive codes and proves this
definition equivalent to reduction from its universal halting predicate;
see \leanmodule{HGAS/Computability/SigmaOne.lean}.

\subsection{The fixed constants}

The constants used in the formal statement are deliberately defined from
the actual compiler, rather than hidden behind an existential quantifier.
Let \(E_0,V_0\) be the fixed row and variable counts of the selected
universal DQL representation.  The repair in Section~\ref{sec:augmentation}
uses
\[
       E_1=E_0+V_0,\qquad V_1=2V_0.
\tag{1.2}
\]
Let \(n_{\rm src}\) be the resulting compact-source dimension and let
\(D_{\rm wrap}\) be its fixed odd wrapper degree.  The homogeneous
product-sphere wrapper uses two source sphere blocks and adds one projective
coordinate, so define
\[
 n_\ast=n_{\rm src}+1,\qquad
 d_\ast=D_{\rm wrap}+2(2+1)=D_{\rm wrap}+6.
\tag{1.3}
\]
The literal \Lean{} definitions are
\leanname{fixedDimension} and \leanname{fixedDegree} in
\leanmodule{HGAS/Arithmetic/NormalizedCorrectedAugmentedEffectiveFamily.lean}.
The same file proves \(n_\ast>0\), \(d_\ast\) odd, and \(d_\ast\geq3\)
as \leanname{fixedDimension_pos}, \leanname{fixedDegree_odd}, and
\leanname{fixedDegree_ge_three}.

Put
\[
 k_\ast=\frac{d_\ast-1}{2},
 \qquad
 N_\ast=\#\{\alpha\in\N^{n_\ast}:|\alpha|=k_\ast\}.
\tag{1.4}
\]
The last cardinality is the dimension of the degree-\(k_\ast\) Veronese
space.  In \Lean{} it is
\[
 \leanname{Veronese.cubicLiftCodeDimension}
   \bigl(\leanname{fixedDimension},
         \leanname{fixedOutputHalfDegree}\bigr).
\]

\begin{theorem}[Fixed-pair and cofinal completeness]
\label{thm:main}
With the fixed constants above:

\begin{enumerate}
\item \(\HGAS_{\Q}(n_\ast,d_\ast)\) is
  \(\Sigma^0_1\)-complete;
\item for every \(N\geq N_\ast\) and every odd \(d\geq3\),
  \(\HGAS_{\Q}(N,d)\) is \(\Sigma^0_1\)-complete; and
\item in particular, for every \(N\geq N_\ast\),
  \(\HGAS_{\Q}(N,3)\) is \(\Sigma^0_1\)-complete.
\end{enumerate}
\end{theorem}

The three kernel-checked declarations are, respectively,
\[
\begin{gathered}
\leanname{fixedPair_sigmaOneComplete},\\
\leanname{cofinal_sigmaOneComplete},\\
\leanname{cubic_sigmaOneComplete}
\end{gathered}
\]
in
\leanmodule{HGAS/Arithmetic/NormalizedCorrectedAugmentedEffectiveFamily.lean}.
The same file records the undecidability consequences as
\leanname{fixedPair_not_computable} and
\leanname{cubic_not_computable}, and the non-c.e.-complement consequences
as \leanname{fixedPair_complement_not_re} and
\leanname{cubic_complement_not_re}.

\begin{remark}[What is not claimed]
The theorem proves cubic undecidability in every sufficiently large fixed
dimension.  It does \emph{not} decide the separate low-dimensional problem
\(\HGAS_{\Q}(3,3)\).  Nothing in the Veronese or padding argument lowers
the fixed dimension to three.
\end{remark}

\begin{remark}[Fixed does not mean numerically printed]
The universal Diophantine representation is a finite object chosen from a
proved existence theorem.  Once chosen, its dimensions and coefficient
tables are constants, and all input-specialization maps are proved
computable.  But the choice is noncomputable in the programming-language
sense, so the formal development does not advertise a command
\texttt{\#eval} that prints \(n_\ast\).  The theorem is an ordinary fixed-pair
existence and many-one completeness theorem, not a numerical
resource-minimization result.
\end{remark}

\subsection{Proof map}

The proof has two branches:
\[
\begin{array}{c}
\text{homogeneous GAS}
 \Longrightarrow
\text{finite exact rational Lyapunov certificate}
 \Longrightarrow
\HGAS_{\Q}(n,d)\in\Sigma^0_1,
\\[5pt]
\text{universal halting}
\Longrightarrow
\text{fixed augmented DQL}
\Longrightarrow
\text{compact source with a strict invariant-average sign}
\\
\Longrightarrow
\text{fixed homogeneous rational field}
\Longrightarrow
\text{cubicization and cofinal padding}.
\end{array}
\tag{1.5}
\]
The remainder of the document proves each arrow and identifies its
\Lean{} counterpart.

\section{The computably enumerable upper bound}
\label{sec:upper}

\begin{theorem}[Uniform c.e. upper bound at fixed odd degree]
\label{thm:upper}
For every fixed \(n\) and odd \(d\), the predicate
\(\HGAS_{\Q}(n,d)\) is computably enumerable.
\end{theorem}

This is \leanname{hgasQ_re_of_oddDegree} in
\leanmodule{HGAS/Lyapunov/RationalHomogeneousConverse.lean}.  We give the
full mathematical argument because the lower bound alone would establish
hardness, not \(\Sigma^0_1\)-completeness.  Its formal scope is one fixed
dimension and one fixed odd degree.  The manuscript's coded
variable-resource and even-degree upper bound instead uses a paper-level
real quantifier-elimination argument and is not a separate final Lean
declaration.

\subsection{The exact certificate}

Write
\[
       \rho(x)=x_1^2+\cdots+x_n^2.
\]
A certificate contains a rational polynomial \(p\), a natural exponent
\(r\), and two natural grid resolutions \(m_p,m_q\).  The checker generates
the finite rational grids, coefficient--derivative bounds, and rational
margins from these data.  The certificate represents the degree-two
homogeneous rational function
\[
             V(x)=\frac{p(x)}{\rho(x)^r}.
\tag{2.1}
\]
If \(F\) has degree \(d\), differentiation gives
\[
\dot V(x)
 =\frac{\rho(x)\nabla p(x)\!\cdot\!F(x)
       -2r\,p(x)\,x\!\cdot\!F(x)}
        {\rho(x)^{r+1}}.
\]
Thus the polynomial numerator used to certify strict decrease is
\[
 q(F,p,r)
   =-\rho\,\nabla p\!\cdot\!F+2r\,p\,x\!\cdot\!F.
\tag{2.2}
\]
The semantic polynomial is
\leanname{HGAS.rationalLyapunovNumerator} in
\leanmodule{HGAS/Lyapunov/RationalCertificateAlgebra.lean}.  Its executable
sparse implementation is
\leanname{PolynomialCode.rationalLyapunovNumerator} in
\leanmodule{HGAS/Foundation/PolynomialDifferentiationCode.lean}; a decoding
theorem identifies the two exactly.

Although \(\rho\) in (2.1) is the squared Euclidean norm, homogeneity makes
strict positivity of \(p\) and \(q\) away from zero equivalent to strict
positivity on the unit sphere of any norm.  The Boolean checker uses the
sup-norm unit sphere, whose boundary admits a particularly simple finite
rational grid.  Exact derivative bounds transfer the checked grid margin
to every point of that sphere.  The finite checker and its semantic
interpretation are developed in
\leanmodule{HGAS/Lyapunov/ExactGridPositivity.lean} and
\leanmodule{HGAS/Lyapunov/GridRationalCertificateSoundness.lean};
primitive recursiveness is proved in
\leanmodule{HGAS/Lyapunov/ExactGridPositivityComputability.lean} and
\leanmodule{HGAS/Lyapunov/GridRationalCertificate.lean}.

If it accepts, (2.1) is positive definite away from zero and (2.2) is
positive away from zero.  Coercivity of \(V\) bounds forward trajectories
inside compact sublevel sets.  Standard local ODE continuation then proves
forward completeness; the strict derivative inequality gives Lyapunov
stability and global attraction.  Hence accepted certificates are sound
for the exact GAS definition, including completeness.  This is proved in
\leanmodule{HGAS/Lyapunov/GridRationalCertificateSoundness.lean}.

\subsection{Completeness of the certificate search}

Suppose now that \(F\) is homogeneous of odd degree \(d\) and is GAS.
The completeness proof has five steps.

\paragraph{1. A homogeneous \(C^1\) Lyapunov function.}
Separate radius and direction.  On the unit sphere put
\[
       a(u)=\langle u,F(u)\rangle,\qquad
       Y(u)=F(u)-a(u)u.
\tag{2.3}
\]
Then \(Y\) is tangent to the sphere.  For \(x=ru\), homogeneity gives,
after the usual angular time change,
\[
       \frac{\dd u}{\dd s}=Y(u),\qquad
       \frac{\dd}{\dd s}\log r=a(u).
\tag{2.4}
\]
GAS yields a uniform negative radial-cocycle block on the compact sphere.
Integrating a suitably smoothed finite-horizon cocycle produces a strict
degree-two homogeneous \(C^1\) Lyapunov function.  The construction is
formalized in
\leanmodule{HGAS/Lyapunov/HomogeneousConverseC1.lean} and
\leanmodule{HGAS/Lyapunov/EuclideanAngularConverse.lean}.  It is closely
related to Rosier's homogeneous converse theorem~\cite{Rosier1992}, but the
particular theorem needed here is proved internally.

\paragraph{2. Evenization.}
Odd degree implies \(F(-x)=-F(x)\).  If \(V\) is a strict Lyapunov
function, then
\[
                 V_{\rm ev}(x)=\frac{V(x)+V(-x)}2
\tag{2.5}
\]
is even, positive, degree two, and has strictly negative orbital
derivative.  The proof, including the derivative identity, is reusable in
\leanmodule{HGAS/Analysis/EvenC1Lyapunov.lean}.

\paragraph{3. Simultaneous \(C^1\) polynomial approximation.}
On a cube containing the unit sphere, tensor Bernstein polynomials
approximate \(V_{\rm ev}\) and all first partial derivatives uniformly.
The development proves the tensor derivative formula rather than assuming
a black-box approximation theorem; see
\leanmodule{HGAS/Analysis/TensorBernsteinPolynomial.lean},
\leanmodule{HGAS/Analysis/TensorBernsteinDerivative.lean}, and
\leanmodule{HGAS/Analysis/TensorBernsteinC1Approximation.lean}.
For sufficiently small \(C^1\) error, positivity and strict decay retain
positive margins on the sphere.

\paragraph{4. Homogenization and rational perturbation.}
Evenness permits reconstruction of an even homogeneous numerator after
multiplication by a sufficiently high power of \(\rho\).  Rational
coefficients are dense in the finite coefficient space, so a sufficiently
small rational perturbation preserves both compact strict margins.  This
produces a rational function of the form (2.1).

\paragraph{5. A finite exact grid witness.}
Strict positivity on a compact sphere supplies positive margins and finite
derivative bounds.  A sufficiently fine rational grid therefore passes the
exact checker.  This is the completeness theorem in
\leanmodule{HGAS/Lyapunov/GridRationalCertificateCompleteness.lean}.

Consequently every positive instance has some finite accepted certificate.
Enumerating all finite candidate codes semidecides GAS.  This rational
certificate viewpoint is consistent with the algebraic Lyapunov results of
Ahmadi and El Khadir~\cite{AhmadiElKhadir2020}; here, however, the exact
certificate language, soundness, approximation, and enumerability chain
are all present in the kernel-checked development.

\section{The universal arithmetic source}
\label{sec:arithmetic}

\subsection{Universal halting}

\Mathlib{} represents partial-recursive functions by a countable type of
codes.  Define
\[
  \mathrm{UniversalHalts}(e,x)
  \quad\Longleftrightarrow\quad
  \operatorname{eval}(e,x)\text{ is defined}.
\tag{3.1}
\]
Its domain is c.e.  Conversely, if \(A\subseteq\N\) is c.e., choose a code
\(e_A\) for the partial computation which returns \(0\) precisely on
members of \(A\).  Then
\[
             x\in A
       \quad\Longleftrightarrow\quad
             \mathrm{UniversalHalts}(e_A,x).
\tag{3.2}
\]
The map \(x\mapsto(e_A,x)\) is computable.  This proves universal-halting
\(\Sigma^0_1\)-completeness in
\leanmodule{HGAS/Computability/SigmaOne.lean}.

For the arithmetic compiler it is convenient to encode the code--input
pair by one natural number.  The predicate
\(\leanname{UniversalHalting.Predicate}:\N\to\mathrm{Prop}\) is the
resulting unary complete predicate.  Its exact encoding and bounded
interpreter are in
\leanmodule{HGAS/Arithmetic/ExecutableMRDPBoundary.lean}.

\subsection{The formal MRDP fragment}

The reduction needs a fixed polynomial \(P(t,z_1,\ldots,z_s)\in
\Z[t,z_1,\ldots,z_s]\) such that
\[
 \mathrm{Predicate}(t)
 \quad\Longleftrightarrow\quad
 \exists z\in\N^s,\ P(t,z)=0.
\tag{3.3}
\]
The project does not introduce (3.3) as an axiom.  It proves the precise
primitive-recursive Diophantine closure statement needed for the bounded
universal-machine trace.

Here is the mathematical core.  Diophantine graphs are plainly closed
under the initial functions and composition.  Primitive recursion asks
one to express a finite trace
\[
 a_0=g(\bar x),\qquad
 a_{i+1}=h(\bar x,i,a_i)\quad(0\leq i<N),\qquad
 a_N=y
\tag{3.4}
\]
with a fixed number of existentially quantified natural variables.  A
sufficiently large radix \(\beta\) packs the sequence into one integer
\[
                     A=\sum_{i=0}^{N}a_i\beta^i.
\tag{3.5}
\]
Sparse radix ciphers provide Diophantine access to the digits and to
finite aligned register operations.  The conditions in (3.4), apparently
a bounded universal quantification over \(i<N\), are combined into a
single sparse bounded-forall certificate.  The radix is itself certified
large enough that carries cannot corrupt any register.  Correctness has
two directions:

\begin{itemize}
\item a genuine finite computation produces the packed certificate; and
\item every accepted packed certificate decodes, digit by digit, to a
  genuine computation satisfying (3.4).
\end{itemize}

The principal modules are
\leanmodule{HGAS/Arithmetic/SparseRadixCipher.lean},
\leanmodule{HGAS/Arithmetic/SparseCipherCircuit.lean},
\leanmodule{HGAS/Arithmetic/SparseBoundedForall.lean}, and
\leanmodule{HGAS/Arithmetic/PrimitiveRecursiveDiophantine.lean}.
Their endpoints include
\leanname{cipherCertificate_iff_boundedForall},
\leanname{boundedForall_dioph},
\leanname{recursionCertificate_iff},
\leanname{natPrimrec_unaryGraph_dioph}, and
\leanname{binaryPrimrecRel_dioph}.

Applying the last theorem to the primitive-recursive bounded interpreter
proves
\[
 \leanname{binaryPrimitiveRecursiveRelationsAreDiophantine}
\]
in \leanmodule{HGAS/Arithmetic/UnconditionalMRDP.lean}.  Finite-support
extraction then supplies (3.3).  This is the portion of the classical
Davis--Putnam--Robinson--Matiyasevich theorem
\cite{DavisPutnamRobinson1961,Matiyasevich1970,Matiyasevich1971,
Matiyasevich1993} needed by the reduction, formalized rather than
postulated.

\subsection{Fixed diagonal--quadratic--linear systems}

A diagonal--quadratic--linear, or DQL, system with \(E\) rows and \(V\)
integer variables is
\[
  \sum_{j=1}^{V}\bigl(A_{ij}x_j^2+B_{ij}x_j\bigr)=c_i,
       \qquad i=1,\ldots,E,
\tag{3.6}
\]
with all coefficients integral.  Mixed quadratic terms are forbidden.
This restriction loses no effectiveness.  A multiplication gate \(xy=z\)
can be replaced by a fresh variable \(w\) and
\[
                  w=x+y,\qquad
                  w^2-x^2-y^2-2z=0.
\tag{3.7}
\]
Natural-number witnesses are represented over \(\Z\) with four squares,
using Lagrange's theorem~\cite{Lagrange1770}.  Addition, constants, copies,
and (3.7) compile a finite polynomial circuit into finitely many DQL rows.
The equivalence of witnesses is proved in
\leanmodule{HGAS/Arithmetic/DQL.lean},
\leanmodule{HGAS/Arithmetic/GateDQLCompiler.lean}, and
\leanmodule{HGAS/Arithmetic/FixedRepresentationDQL.lean}.

After the polynomial in (3.3) has been fixed, substituting the input
\(t\) changes only right-hand-side data in a fixed template.  Thus there
are fixed \(E_0,V_0\) and a computable map
\[
       t\longmapsto I_t\in\mathrm{DQL}(E_0,V_0)
\tag{3.8}
\]
such that
\[
             I_t\text{ is solvable over }\Z
       \quad\Longleftrightarrow\quad
             \mathrm{Predicate}(t).
\tag{3.9}
\]
These are
\leanname{unconditionalUniversalHaltingFixedDQL},
\leanname{unconditionalUniversalHaltingFixedDQL_computable}, and
\leanname{unconditionalUniversalHaltingFixedDQL_solvable_iff} in
\leanmodule{HGAS/Arithmetic/UnconditionalFixedDQL.lean}.

\subsection{The nonzero-linear-column augmentation}
\label{sec:augmentation}

The final three-moment selector needs the linear character of every
arithmetic variable to be nonzero.  A general DQL representation need not
have this property: a variable may occur only through \(x_j^2\).  We enforce
it without changing solvability.

For every old variable \(x_j\), add a duplicate \(y_j\) and the row
\[
                        x_j-y_j=0.
\tag{3.10}
\]
Old rows use only the \(x\)-variables.  If \(x\) solves the old system,
\((x,x)\) solves the augmented one.  Conversely, the first half of every
augmented solution solves all old rows.  Therefore
\[
                 \operatorname{Aug}(I)\text{ solvable}
            \quad\Longleftrightarrow\quad
                 I\text{ solvable}.
\tag{3.11}
\]
Every old column contains coefficient \(+1\) in its new equality row, and
every duplicate column contains coefficient \(-1\).  Hence every augmented
linear column is nonzero.

The construction and both implications in (3.11) are in
\leanmodule{HGAS/Arithmetic/DQLNonzeroLinearAugmentation.lean}.  The same
file proves that augmentation is primitive recursive.  Applied to (3.8),
it gives fixed dimensions (1.2), a computable family
\(\widetilde I_t\), and
\[
\begin{split}
\widetilde I_t\text{ solvable}
 &\Longleftrightarrow \mathrm{Predicate}(t),\\
\forall j\ \exists i\quad
   \widetilde B_{ij}&\neq0.
\end{split}
\tag{3.12}
\]
The fixed-resource specialization is in
\leanmodule{HGAS/Arithmetic/NormalizedCorrectedAugmentedUniversalResources.lean}.

\section{Homogeneous GAS as a compact invariant-average sign}
\label{sec:angular-sign}

The compiler first constructs a compact polynomial source and a continuous
reward.  This section explains why the sign of the largest invariant
average is exactly the right interface to homogeneous GAS.

\subsection{Polar dynamics}

Suppose \(F\) is homogeneous of degree \(m+1\), where \(m>0\).  For a unit
vector \(u\), define \(a,Y\) by (2.3).  If a nonzero trajectory is written
\(x(t)=r(t)u(t)\), then
\[
       \dot r=r^{m+1}a(u),\qquad
       \dot u=r^mY(u).
\tag{4.1}
\]
With angular time \(\dd s/\dd t=r^m\),
\[
       u'=Y(u),\qquad
       (\log r)'=a(u).
\tag{4.2}
\]
Consequently,
\[
       \log\frac{r(s)}{r(0)}
          =\int_0^s a(\phi_\tau u(0))\,\dd\tau,
\tag{4.3}
\]
where \(\phi\) is the angular flow on the compact sphere.

\subsection{Uniform negative blocks}

For a compact flow \(\phi\) and continuous reward \(a\), write
\[
       \mathcal M_\phi(a)
       =\max_{\mu\in\mathcal I(\phi)}\int a\,\dd\mu,
\tag{4.4}
\]
where \(\mathcal I(\phi)\) is the set of invariant Borel probabilities.
The maximum exists.  The following equivalence is the key compactness
principle:
\[
\mathcal M_\phi(a)<0
\quad\Longleftrightarrow\quad
\exists T,c>0\ \forall x,\quad
   \int_0^T a(\phi_t x)\,\dd t\leq-c.
\tag{4.5}
\]

For the forward implication, if no uniform block exists, choose longer
orbit segments with average arbitrarily close to zero.  Their occupation
probabilities have a weakly convergent subnet.  The endpoint error in
invariance is \(O(1/T)\), so the limit is invariant and has nonnegative
reward average, contradicting (4.4).  For the reverse implication,
integrate the uniform block against any invariant probability.

The occupation-measure argument, existence and compactness of invariant
probabilities, and the equivalence (4.5) are formalized in
\leanmodule{HGAS/Dynamics/KrylovBogolyubov.lean},
\leanmodule{HGAS/Dynamics/InvariantProbabilityCompact.lean}, and
\leanmodule{HGAS/Dynamics/UniformNegativeBlocks.lean}.

Combining (4.2)--(4.5), including the conversion back to physical time,
gives
\[
       0\text{ is GAS for }F
       \quad\Longleftrightarrow\quad
       \mathcal M_\phi(a)<0.
\tag{4.6}
\]
The reusable formal endpoint is
\leanname{globallyAsymptoticallyStable_iff_maximalAngularAverage_neg} in
\leanmodule{HGAS/Wrapper/AngularAverageGAS.lean}.

\subsection{An ergodic maximizer}

The compact-accumulator argument needs a maximizing invariant probability
which is ergodic for the entire real flow, not merely for its time-one map.
The project proves directly that the compact convex set of invariant
probabilities has an extreme maximizer for the affine functional
\(\mu\mapsto\int a\,\dd\mu\).  It embeds measures into the weak dual of
bounded continuous functions, applies Krein--Milman to the maximizing
face, and proves that an extreme invariant probability satisfies the
full-flow zero--one law by splitting on an invariant measurable set.

The result is
\leanname{exists_ergodicInvariantProbability_isMaxOn_average} in
\leanmodule{HGAS/Dynamics/ErgodicMaximizer.lean}.  In particular, to upper
bound (4.4), it suffices to upper bound the reward average of every
full-flow ergodic invariant probability.  No measurable ergodic
decomposition theorem is used.

\section{The arithmetic Fourier detector}
\label{sec:detector}

Let \(I=(A,B,c)\) be a DQL system with \(V\) variables and \(E\) rows.
Use a torus
\[
           \T^{V+E}
       =(\R/\Z)^{\{\text{private variables}\}\sqcup
                    \{\text{equation rows}\}}.
\]
For \(x,s\in\Z^V\), define the private and equation exponents
\[
\begin{split}
 p_j(x,s)&=x_j^2-s_j,\\
 q_i(x,s)&=\sum_j(A_{ij}s_j+B_{ij}x_j)-c_i.
\end{split}
\tag{5.1}
\]
The full exponent vector is zero exactly when
\[
             s_j=x_j^2\ \text{for all }j,
       \qquad I(x)=0.
\tag{5.2}
\]

Choose \(0<\lambda,\eta<1\).  The expanded detector is the absolutely
convergent Fourier series
\[
 \mathcal D_I(\theta)=
 \sum_{x,s\in\Z^V}
   \lambda^{\sum_j|x_j|^2}
   \eta^{\sum_j|s_j|}
   e^{2\pi i\langle (p(x,s),q(x,s)),\theta\rangle}.
\tag{5.3}
\]
Haar integration deletes every nonzero Fourier character.  By (5.2),
\[
 \int_{\T^{V+E}}\mathcal D_I\,\dd\theta
  =
 \sum_{\substack{x\in\Z^V\\I(x)=0}}
       (\lambda\eta)^{\sum_j|x_j|^2}.
\tag{5.4}
\]
Every surviving term is strictly positive.  Therefore
\[
 \int\mathcal D_I\,\dd\theta\neq0
       \quad\Longleftrightarrow\quad
 I\text{ is solvable}.
\tag{5.5}
\]
There is no cancellation argument hidden here: the right side of (5.4) is
a nonnegative real series.

The exact support calculation is in
\leanmodule{HGAS/Arithmetic/PositiveConstantTerm.lean}; absolute
summability, continuity, Haar extraction, and (5.5) are in
\leanmodule{HGAS/Arithmetic/DetectorFourier.lean}.  The final compiler uses
the concrete values
\[
                         \lambda=\frac1{256},
            \qquad       \eta=\frac12.
\tag{5.6}
\]

\subsection{Prime-radical clocks}

Assign a distinct prime \(p_\ell\) to every fundamental phase coordinate.
A polynomial controller on a unit circle has locked equilibria satisfying
\[
          p_\ell r_\ell^2=1,
\qquad
          \omega_\ell=p_\ell r_\ell
              =\varepsilon_\ell\sqrt{p_\ell},
\qquad
          \varepsilon_\ell\in\{+1,-1\}.
\tag{5.7}
\]
The phase circle rotates at angular velocity \(\omega_\ell\).  Under
\(\R/\Z\to S^1,\ \theta\mapsto e^{2\pi i\theta}\), its frequency in turns
is \(\omega_\ell/(2\pi)\).  The detector uses fourth powers of the
fundamental phases, so its actual frequency is
\[
               \Omega_\ell
             =\frac{2}{\pi}\,
                \varepsilon_\ell\sqrt{p_\ell}.
\tag{5.8}
\]

Square roots of distinct primes are linearly independent over \(\Q\).
The project proves this by adjoining the roots one at a time and using
the square-class obstruction; it does not import the result as an axiom.
Signs and the common nonzero scale \(2/\pi\) preserve independence.
Hence
\[
      \sum_\ell k_\ell\Omega_\ell=0,\quad k_\ell\in\Z
      \quad\Longrightarrow\quad k_\ell=0\text{ for all }\ell.
\tag{5.9}
\]
See \leanmodule{HGAS/Clocks/RadicalIndependence.lean},
\leanmodule{HGAS/Clocks/ScaledRadicalClock.lean}, and
\leanmodule{HGAS/Arithmetic/NormalizedCorrectedEffectiveSelectedDetectorPhase.lean}.

For any initial phase and any horizon sequence \(T_N\to\infty\), termwise
Cesàro integration followed by dominated summation gives
\[
 \frac1{T_N}\int_0^{T_N}
        \mathcal D_I(\theta_0+t\Omega)\,\dd t
       \longrightarrow
 \int_{\T^{V+E}}\mathcal D_I\,\dd\theta.
\tag{5.10}
\]
The arbitrary-horizon form is important because the recurrent orbit
selected later supplies its own physical return horizons.  It is proved in
\leanmodule{HGAS/Analysis/TorusFourierCesaroArbitraryHorizon.lean} and
\leanmodule{HGAS/Arithmetic/DetectorFourierCesaroArbitraryHorizon.lean}.

\section{The polynomial modular--Jacobi source}
\label{sec:jacobi-source}

The detector (5.3) is an analytic Fourier series, whereas the reduction
must emit a finite polynomial vector field with rational coefficients.
The modular--Jacobi source realizes the needed theta factors as coordinates
of an autonomous polynomial ODE.

\subsection{Unary and Jacobi coordinates}

The unary modular block has four complex coordinates
\[
                     (A,C,E,Y).
\]
On the canonical theta solution these are, up to the fixed normalization,
\[
 A=\vartheta_2^4,\qquad
 C=\vartheta_4^4,\qquad
 Y=\vartheta_3,\qquad
 E=E_2.
\tag{6.1}
\]
Ramanujan/Jacobi differential identities make their derivative under
horizontal motion a polynomial in \((A,C,E,Y)\).

The reciprocal Jacobi block has ten complex coordinates
\[
 (p_2,p_3,p_4,H,T_1,T_2,T_3,T_4,\Delta,U).
\tag{6.2}
\]
The first four encode theta-square and Eisenstein coefficient data;
\(T_1,\ldots,T_4,\Delta\) encode an elliptic theta state; and
\[
                          U\,T_1=1
\tag{6.3}
\]
keeps the trajectory in a regular reciprocal chart.  There are two
commuting rational polynomial derivations: an elliptic derivation and a
modular derivation.  Their literal sparse Gaussian-rational codes and
realification to twenty real coordinates are in
\leanmodule{HGAS/Modular/JacobiPolynomialCode.lean}.

The unary and Jacobi blocks are linked, but retained as separate ambient
coordinates.  Ten polynomial residuals enforce the coefficient identities,
the two theta quadrics, the reciprocal chart, and a coherent signed root.
These residuals form an invariant ideal.  The combined state therefore has
fourteen complex, or twenty-eight real, coordinates and an exact finite
polynomial compiler; see
\leanmodule{HGAS/Modular/LinkedJacobiPolynomialCode.lean}.
The formulas are related to the polynomial theta systems studied by
Brezhnev, Lyakhovich, and Sharapov~\cite{BrezhnevLyakhovichSharapov2011}
and Brezhnev~\cite{Brezhnev2013}.  All normalizations actually used in the
reduction are rederived and checked in the project.

\subsection{Controllers, phase factors, and residual drifts}

There is one linked unary--Jacobi column for each DQL variable.  Polynomial
controller circles provide the two clock coefficients of each column.
Additional real polynomial coordinates provide:

\begin{itemize}
\item the fourth-root gauge multiplying \(T_1\);
\item two corrected first-moment selectors;
\item reciprocal Poisson factors with parameter \(1/2\);
\item a right-hand-side phase oscillator; and
\item nonnegative validity and clock-radius residual observables.
\end{itemize}

Complex polynomial codes are realified mechanically.  The controller and
phase circle equations are tangent to their unit circles at every ambient
point; the linked residual ideal is preserved; and every displayed ODE
identity is proved by evaluation of the finite sparse code.  These pieces
are assembled in the sequence beginning with
\leanmodule{HGAS/Arithmetic/FixedDetectorPolynomialCode.lean} and ending
with
\leanmodule{HGAS/Arithmetic/NormalizedCorrectedFixedDetectorCompactSourceCode.lean}.

The word ``normalized'' records three important corrections which are
literal in the code:

\begin{enumerate}
\item angular velocities are converted from radians to turns by \(2\pi\);
\item detector phases use fourth powers, hence the factor \(4\); and
\item the two first moments are centered at the compiler's exact constant
  \(1/256\).
\end{enumerate}

\subsection{From zero conditional drift to a selected orbit}

Later, the compact accumulator will imply that every drift coordinate has
zero integral under a certain conditional invariant probability.  Apply
the pointwise ergodic theorem simultaneously to the finite family of drift,
speed, and validity coordinates.  Poincar\'e recurrence for the flow and
its reverse supplies a point which:

\begin{itemize}
\item lies in the finite central chart;
\item is recurrent both forward and backward in compact time;
\item has zero validity energy, so all circle and linked residual equations
  hold pointwise;
\item realizes the prescribed zero drift means; and
\item has a positive average physical-speed coordinate.
\end{itemize}

The selected object is
\leanname{SelectedBidirectionalValidityOrbit}.  Its construction from
conditional zero means is
\leanname{exists_selectedBidirectionalValidityOrbit_of_conditionalBase_zero}
in
\leanmodule{HGAS/Arithmetic/NormalizedCorrectedFixedDetectorBidirectionalValidityOrbit.lean}.
The bidirectional clause is not cosmetic: standardizing a column by a
negative clock frequency reverses time, so forward recurrence alone would
not suffice.

\section{Classifying a selected Jacobi column}
\label{sec:classification}

Fix one selected detector column \(j\).  Write \(\omega_j\neq0\) for its
private modular frequency and \(N_j\) for its linear DQL frequency.  Rescale
time by
\[
                         a_j=\frac{\pi}{\omega_j}.
\tag{7.1}
\]
The standardized modular clock is exactly \(\pi i\), while its real
elliptic frequency \(\nu_j\) is defined by
\[
             \pi\nu_j
       =a_j\,\frac{\omega_j+N_j}{2}.
\tag{7.2}
\]
Both identities, with all signs retained, are proved in
\leanmodule{HGAS/Arithmetic/NormalizedCorrectedFixedDetectorColumnStandardization.lean}
and
\leanmodule{HGAS/Arithmetic/NormalizedCorrectedFixedDetectorCanonicalRootSeed.lean}.

\subsection{The level-two horocycle alternative}

The unary coefficient curve determines a canonical modular
\(\lambda\)-coordinate.  The project reconstructs a lift to the upper
half-plane and an exact real-frame horocycle.  On the level-two frame
quotient, a nonperiodic horocycle is equidistributed for normalized Haar
probability, whereas periodic horocycles are closed cusp orbits.  The
specialized Dani--Smillie theorem is proved through:

\begin{enumerate}
\item quantitative nonescape of occupation measures;
\item avoidance of mass on the countable periodic family;
\item good return times for compactly supported tests;
\item the nonergodic pointwise ergodic theorem; and
\item uniqueness of an invariant probability off the periodic set.
\end{enumerate}

The final unconditional statements are
\leanname{gammaTwoPointwiseHorocycleEquidistribution_unconditional} and
\leanname{gammaTwoDaniSmillieTheorem_unconditional} in
\leanmodule{HGAS/Modular/GammaTwoOffPeriodicUnique.lean}.  The argument
formalizes the level-two instance of Dani's invariant-probability
classification~\cite{Dani1978} and the periodic-or-equidistributed orbit
dichotomy of Dani and Smillie~\cite{DaniSmillie1984}; Hubbard and Miller
provide a modern finite-area treatment of the latter
\cite{HubbardMiller2010}.

The selected detector orbit cannot lie in the Haar-generic branch.  A
compactly supported family of theta tests has finite value along the
selected compact orbit, but its Haar cusp moment diverges.  Exact
orientation reversal and constant time rescaling handle either sign of
\(\omega_j\).  Therefore the reconstructed frame orbit is periodic.
The composition is
\leanname{standardizedUnaryColumn_hasHorizontalOrPeriodicExactFrameLift_unconditional}
in
\leanmodule{HGAS/Arithmetic/NormalizedCorrectedFixedDetectorStandardizedExactFrameExclusion.lean}.

\subsection{The finite closed-cusp selector}

A periodic frame orbit admits a horizontal chart
\[
                    \tau(t)=x+s t+i y,
             \qquad y>0,\quad s\neq0.
\tag{7.3}
\]
The two compiled raw moments have Cesàro limit \(1\).  Fourier series for
\(\vartheta_2,\vartheta_3,\vartheta_4\) on (7.3) show that the wrong cusp
branches have zero mean.  In the surviving infinity cusp, the means force
the coherent root scales and quartic amplitude to equal \(1\).  Thus,
after the exact resampling back from (7.1),
\[
             Y(t)=\vartheta_3(\tau(t)),\qquad
             p_4(t)=\vartheta_4(\tau(t))^2.
\tag{7.4}
\]
This is proved in
\leanmodule{HGAS/Modular/FiniteClosedCuspMeanSelector.lean},
\leanmodule{HGAS/Modular/FiniteClosedCuspResampledMeanSelector.lean}, and
\leanmodule{HGAS/Arithmetic/NormalizedCorrectedFixedDetectorFiniteClosedBranchSelection.lean}.

The invariant residuals now recover every unary coefficient:
\[
\begin{split}
 p_4^2-C=0&\Longrightarrow C=\vartheta_4^4,\\
 Y^4-A-C=0&\Longrightarrow A=\vartheta_2^4,\\
 12H-E=0&\Longrightarrow H=E_2/12.
\end{split}
\tag{7.5}
\]
For the last identity, the common nonvanishing \(Y=\vartheta_3\) coordinate
and its polynomial differential equation allow cancellation of the clock
and \(Y\), forcing equality of the two \(E\)-coordinates.

Only the sign of \(p_2\) remains.  Since
\[
             p_2^2=\vartheta_2(\tau)^4,
\]
we have \(p_2=\pm\vartheta_2(\tau)^2\).  Translation
\(\tau\mapsto\tau+2\) changes the sign of \(\vartheta_2^2\) while fixing
the other canonical root coordinates.  Hence one of \(\tau\) and
\(\tau+2\) is an exact canonical seed for
\((p_2,p_3,p_4,H)\).  This argument is
\leanname{exists_canonicalRootSeed_of_finiteClosedSelection} in
\leanmodule{HGAS/Modular/FiniteClosedCuspCanonicalRootSeed.lean}.

\subsection{Global affine Bloch continuation}

At the canonical seed, the two Jacobi quadrics vanish.  Along the complete
trajectory, (6.3) prevents the projective coordinate \(T_1\) from
vanishing.  Pointwise theta saturation reconstructs affine Bloch
parameters, and uniqueness for the regular polynomial ODE propagates the
identification for all real time.  Thus the physical column is a globally
regular affine Jacobi--Bloch trajectory.

The detector-facing result is
\leanname{physicalColumn_exists_classifiedSeed} in
\leanmodule{HGAS/Arithmetic/NormalizedCorrectedFixedDetectorPhysicalClassifiedSeed.lean}.
It also chooses the actual initial phase and the unique free gauge constant
matching the physical state at time zero.  No canonical all-one phase is
substituted.

\section{The repaired three-moment theorem}
\label{sec:three-moment}

This is the most delicate identification in the lower bound.  We give it
in detail, including the failed zero-ratio branch which motivated the DQL
augmentation.

\subsection{Quadratic--geometric coefficients}

After multiplication by the fourth-root gauge, an affine Bloch theta trace
has coefficients
\[
                   C_k=C\,R^{k^2}S^k,
                  \qquad k\in\Z,
\tag{8.1}
\]
with \(R,S\neq0\).  Integer translations \(k\mapsto k-\ell\) and reflected
translations \(k\mapsto-k-\ell\) retain the same form.  The direct
factorization is proved in
\leanmodule{HGAS/Modular/AffineBlochQuadraticGeometric.lean}.

A normalized sequence \(R^{k^2}S^k\) is determined by its coefficients at
\(k=1\) and \(k=-1\).  Indeed their product determines \(R^2\), so either
\((R,S)=(R',S')\) or \((R,S)=(-R',-S')\).  The second alternative gives the
same sequence because \(k^2+k\) is even.  This elementary rigidity lemma is
formalized in
\leanmodule{HGAS/Modular/QuadraticGeometricSequence.lean}.

\subsection{What the three moments select}

Let \(\Omega\neq0\) be the private angular frequency and put
\[
                       \rho=\frac{N}{\Omega}.
\tag{8.2}
\]
The quadratic dispersion of the trace makes its mean resonance equation
factor into two Bloch classes.  If \(\rho\) is irrational, the two possible
classes are disjoint:
\[
        a\in\Z
        \qquad\hbox{or}\qquad
        a\in-\rho+\Z.
\tag{8.3}
\]
In the first class, the zeroth and the two twisted first moments select
\[
              C_{-\ell},\quad C_{1-\ell},\quad C_{-1-\ell}.
\tag{8.4}
\]
In the reflected class they select the same three neighboring coefficients
in reverse order.  The normalized zeroth moment is \(1\); the \(+1\) and
\(-1\) moments match the corresponding target theta coefficients.  By the
quadratic--geometric lemma, every coefficient matches the target, hence the
whole theta trace matches.

The twisting constants are the actual initial clock characters.  Their
norm is one, so they are nonzero and may be cancelled; they need not equal
one.  Arbitrary recurrent horizon sequences are permitted.  The exact
phase-free theorem is
\leanname{quadraticTrace_eq_target_of_three_moments} in
\leanmodule{HGAS/Modular/JacobiThreeMomentTargetRigidity.lean}.

\subsection{Why \(\rho=0\) cannot be ignored}

\begin{warning}[A genuine discarded argument]
When \(\rho=0\), the direct and reflected resonance classes in (8.3)
coincide.  The three averages may then see sums of two coefficients rather
than one coefficient, and the preceding proof is false.
\end{warning}

This is not merely a missing side condition.  There are
quadratic--geometric counterfamilies
\[
                 C_k=\lambda^{k^2}A^{k^2}B^k,
       \qquad |A|=1,\qquad B+B^{-1}=2A^{-1},
\tag{8.5}
\]
whose zeroth and two first resonant moments have the prescribed values
although the trace is not the canonical target.  For example, one may take
\[
                    A=i,\qquad
                    B=i(\sqrt2-1).
\tag{8.6}
\]
Thus an older theorem which admitted the zero-ratio branch is not used in
the final dependency chain.

The repair is exactly Section~\ref{sec:augmentation}.  Let
\(\ell_j\) be the nonzero linear DQL character of column \(j\), and let
\(e_j\) be its private character.  The selected torus frequency
\(\Omega_{\rm tor}\) has no nontrivial integer annihilator.  If
\[
       \frac{\langle\ell_j,\Omega_{\rm tor}\rangle}
            {\langle e_j,\Omega_{\rm tor}\rangle}
          =\frac ab\in\Q,
\tag{8.7}
\]
then
\[
             b\ell_j-ae_j
\tag{8.8}
\]
annihilates \(\Omega_{\rm tor}\).  Nonresonance makes (8.8) the zero
character.  Its private coordinate gives \(a=0\); a nonzero equation
coordinate of \(\ell_j\) then gives \(b=0\), contradicting the choice of
denominator.  Therefore (8.7) is irrational for every augmented column.

This finite-character proof is
\leanname{linear_div_private_pairing_irrational} and its physical-frequency
specialization
\leanname{numberFrequency_div_omegaFrequency_irrational} in
\leanmodule{HGAS/Arithmetic/NormalizedCorrectedEffectiveColumnFrequencyIrrational.lean}.

\subsection{Canonical theta output}

Freeze the target Gaussian theta coefficients at the actual initial
detector phase.  They are quadratic--geometric, and the compiled \(+1\)
and \(-1\) selector characters multiply the two neighboring target
coefficients to \(1/256\).  The corrected moment limits and irrationality
argument therefore give, for every column \(j\) and every \(t\geq0\),
\[
       \mathrm{gauge}_j(t)\,T_{1,j}(t)
       =
       \sum_{n\in\Z}
         \mathrm{thetaTerm}_{I,j,n}
           \bigl(1/256,\theta(t)\bigr).
\tag{8.9}
\]
This is exactly
\leanname{hasCanonicalThetaOutputs} in
\leanmodule{HGAS/Arithmetic/NormalizedCorrectedEffectiveCanonicalThetaOutputs.lean}.
Its sole structural premise is that every linear column is nonzero, which
(3.12) supplies for every final universal instance.

\section{The yes-instance nonzero-mean theorem}
\label{sec:yes}

Assume \(I\) is solvable and suppose, toward a contradiction, that every
conditional compiled drift coordinate has mean zero.  Section
\ref{sec:jacobi-source} then supplies a selected bidirectionally recurrent
validity orbit \(O\).

\subsection{Exact factorization of the physical detector}

The selected orbit begins at an arbitrary physical phase.  Let
\(\theta(t)=\theta(0)+t\Omega\) be its detector phase in turns.  The
right-hand-side oscillator has the same frequency as the flattened DQL
right-hand-side character, including the compiler's factor \(-4\).
Uniqueness for the corresponding linear complex ODE gives
\[
          R_{\rm rhs}(t)
       = \zeta\,
          \chi_{\rm rhs}(\theta(t)),
\qquad
          \zeta\neq0.
\tag{9.1}
\]
The constant \(\zeta\) is retained; it is not silently normalized to one.

The reciprocal graph residuals give the two Poisson factors
\[
       G_+=(1-\tfrac12 W)^{-1},\qquad
       G_-=(1-\tfrac12 W^{-1})^{-1},
\tag{9.2}
\]
where \(|W|=1\).  The absolutely convergent two-sided geometric identity
\[
 \sum_{s\in\Z}\left(\frac12\right)^{|s|}W^s
  =
 \frac{3/4}{(1-\tfrac12W)(1-\tfrac12W^{-1})}
\tag{9.3}
\]
identifies the compiler's literal factor \((3/4)G_+G_-\) with the desired
Poisson series.

Substituting the canonical theta output (8.9), multiplying the finitely
many column factors, and using absolute summability yields the pointwise
identity
\[
          \mathrm{detectorValue}_I(O(t))
       =  \zeta\,\mathcal D_I(\theta(t))
\qquad(t\geq0).
\tag{9.4}
\]
All phase, sign, fourth-power, \(2\pi\), \(3/4\), and \(1/256\)
normalizations in (9.1)--(9.4) are proved in
\leanmodule{HGAS/Arithmetic/NormalizedCorrectedFixedDetectorSelectedPhysicalPhase.lean}
and
\leanmodule{HGAS/Arithmetic/NormalizedCorrectedFixedDetectorSelectedFactorization.lean}.

\subsection{Contradictory Ces\`aro limits}

By (5.9), (5.10), and (9.4), along the selected physical horizons,
\[
 \frac1{T_N}\int_0^{T_N}
       \mathrm{detectorValue}_I(O(t))\,\dd t
   \longrightarrow
 \zeta\int_{\T^{V+E}}\mathcal D_I\,\dd\theta.
\tag{9.5}
\]
Because \(I\) is solvable, (5.5) makes the integral nonzero; because
\(\zeta\neq0\), the limit in (9.5) is nonzero.

On the other hand, the selected orbit was obtained from the assumption
that every conditional drift mean is zero.  The detector drift coordinate
therefore has Ces\`aro limit zero along the same physical horizons.  A
Hausdorff space has unique limits, giving a contradiction.

The analytic statement is
\leanname{false_of_solvable_and_canonicalThetaOutputs} in
\leanmodule{HGAS/Arithmetic/NormalizedCorrectedEffectiveSelectedDetectorAverage.lean}.
The passage back to conditional invariant means is
\leanname{not_conditionalBaseDriftZero_of_solvable} in
\leanmodule{HGAS/Arithmetic/NormalizedCorrectedEffectiveNonzeroMeanFromTheta.lean}.
Consequently:
if \(\mu\) is an ergodic invariant probability whose finite chart has
positive mass, and \(\mu_{\rm fin}\) is its conditional probability there,
then
\[
 I\text{ solvable}
 \quad\Longrightarrow\quad
 \exists q\qquad
       \int b_{I,q}\,\dd\mu_{\rm fin}\ne0 .
\tag{9.6}
\]

\section{The no-instance invariant finite graph}
\label{sec:no}

Now assume \(I\) is not solvable.  The construction must produce one
invariant probability with strictly positive reward.  This branch does not
use the selected-orbit contradiction from Section~\ref{sec:yes}.

\subsection{The orbit-closure Fourier model}

Choose the explicit positive unit equilibrium of every prime controller.
Its phase flow is a translation with frequencies
\[
                 \frac1{2\pi}
                  \bigl(\varepsilon_\ell\sqrt{p_\ell}\bigr)_\ell.
\tag{10.1}
\]
The closure of its orbit is a compact torus, and the orbit is dense there.
Every compiled drift coordinate restricts to an absolutely convergent real
Fourier family
\[
                  g_q(\theta)
          =\sum_{r\in\mathcal I_q}
                c_{q,r}\,
                e^{2\pi i\langle k_{q,r},\theta\rangle}.
\tag{10.2}
\]
The complete finite-code-to-Fourier crosswalk is in
\leanmodule{HGAS/Analysis/NormalizedCorrectedCanonicalDriftFourierFamily.lean}
and
\leanmodule{HGAS/Arithmetic/NormalizedCorrectedCanonicalOrbitClosureModel.lean}.

When \(I\) has no solution, every zero Fourier mode in (10.2) vanishes.
For the detector coordinate this is precisely the contrapositive of the
constant-term calculation (5.4); the other validity and centered-moment
coordinates have zero modes by their explicit normalization.  Thus
\[
                         c_{q,r}=0
                    \quad\hbox{whenever}\quad
                         k_{q,r}=0.
\tag{10.3}
\]

\subsection{Small divisors and a bounded primitive}

For a nonzero integer character \(k\), its frequency
\[
            \langle k,\Omega\rangle
       =\frac1{2\pi}\sum_\ell
              k_\ell\varepsilon_\ell\sqrt{p_\ell}
\tag{10.4}
\]
is nonzero.  More is needed: division by (10.4) must preserve convergence.
The sum in (10.4) is an algebraic number in the finite multiquadratic
extension generated by the selected square roots.  Taking the product of
all Galois conjugates gives a nonzero rational integer after denominators
are cleared.  Bounding all conjugate factors except the original one gives
an effective lower bound for
\(|\langle k,\Omega\rangle|\) in terms of the radical conjugate height of
\(k\).

The detector coefficients have Gaussian/geometric decay.  The development
proves the weighted summability
\[
  \sum_r |c_{q,r}|\,
     H(k_{q,r})^{D_{\rm rad}-1}<\infty,
\tag{10.5}
\]
which dominates the reciprocal small divisor.  Therefore
\[
 U_q(\theta)
   =\sum_{r:k_{q,r}\neq0}
       \frac{c_{q,r}}
            {2\pi i\langle k_{q,r},\Omega\rangle}
       e^{2\pi i\langle k_{q,r},\theta\rangle}
\tag{10.6}
\]
converges uniformly and differentiates along the translation flow to
\(g_q\).  The radical lower bound, summability budget, and primitive
construction occur in
\leanmodule{HGAS/Clocks/RadicalSmallDivisor.lean},
\leanmodule{HGAS/Analysis/CanonicalEffectivePrimitiveBudget.lean}, and
\leanmodule{HGAS/Analysis/NormalizedCorrectedCanonicalEffectivePrimitiveBudget.lean}.

\subsection{Embedding the primitive as an invariant graph}

Scale the primitive vector by a positive rational gain \(\kappa\).  The
gain is chosen effectively from exact finite coefficient bounds so that
the graph lies well inside the finite stereographic accumulator chart.
The base modular/Jacobi orbit is simultaneously centrally compactified.
There is a fixed constant \(c_{\rm cen}>0\) such that the product \(W_0\)
of squared central chart weights satisfies
\[
                       W_0\geq c_{\rm cen}
\tag{10.7}
\]
on the selected orbit closure.  The gain is small enough that the
stereographic scalar coordinate \(W\) of the primitive graph satisfies
\[
                       W\geq1-\frac{c_{\rm cen}}4.
\tag{10.8}
\]

The source reward is
\[
                   \mathcal R=W+2W_0-1-\varepsilon.
\tag{10.9}
\]
The effective reward margin satisfies
\(\varepsilon\leq3c_{\rm cen}/8\).  The formal estimates give
\[
                   \mathcal R\geq
                   \delta:=\frac{3c_{\rm cen}}8>0
\tag{10.10}
\]
on the invariant graph.  Pushing Haar probability on the orbit-closure
torus to this graph produces the required invariant probability.

The exact endpoint is \leanname{finite_graph} in
\leanmodule{HGAS/Arithmetic/NormalizedCorrectedEffectiveNoInstanceFiniteGraph.lean};
its explicit margin is proved in
\leanmodule{HGAS/Arithmetic/NormalizedCorrectedFixedDetectorExplicitFiniteGraphGap.lean}.

\section{Compactification and the strict sign dichotomy}
\label{sec:compact-sign}

\subsection{Central compactification}

Affine source blocks are mapped to central spheres.  If \(y=x/w\) is an
affine chart and \(P(y)\) is polynomial, multiplying by a sufficiently
large common even power of the central weights clears every denominator.
The resulting polynomial field is tangent to every sphere factor,
including points on the equator.  Inside the finite chart it is a positive
time change of the original source.  The sparse rational compiler and its
global tangency theorem are in
\leanmodule{HGAS/Compactification/MixedCentralPolynomialCode.lean},
\leanmodule{HGAS/Compactification/MixedCentralPolynomialField.lean}, and
\leanmodule{HGAS/Compactification/MixedCentralPolynomialGlobalTangency.lean}.

The final normalized source uses two sphere blocks: the centrally
compactified base and the accumulator sphere.  Degree bounds and the
common clearing exponent are computed from finite syntax and are constant
throughout the fixed universal family.

\subsection{The stereographic accumulator}

For a drift vector \(g=(g_1,\ldots,g_m)\), adjoin a sphere coordinate
\((U,W)\in\R^m\times\R\), \(\|U\|^2+W^2=1\), with equations
\[
\begin{split}
 U_i'&=(1+W)g_i-(U\cdot g)U_i,\\
 W'&=-(1+W)(U\cdot g).
\end{split}
\tag{11.1}
\]
These are the inverse-stereographic compactification of translation by
the primitive of \(g\).  Direct differentiation proves tangency to the
accumulator sphere.  The literal polynomial code and evaluation theorem
are in \leanmodule{HGAS/Accumulator/PolynomialCode.lean}.

If an invariant probability gives positive mass to the finite chart
\(W\neq-1\), condition it on that invariant chart and integrate the
derivative of the inverse stereographic coordinates.  Invariance makes
each derivative integral zero; the chart equations reduce this to
\[
                           \int g_i\,\dd\mu=0
\quad\hbox{for every }i.
\tag{11.2}
\]
This compact drift obstruction is proved in
\leanmodule{HGAS/Accumulator/CompactDriftObstruction.lean}.

\subsection{Reward on boundary and infinity}

Recall (10.9).  On a central boundary face, \(W_0=0\), and \(W\leq1\), so
\[
                            \mathcal R\leq-\varepsilon.
\tag{11.3}
\]
At accumulator infinity, \(W=-1\) and \(W_0\leq1\), giving the same bound.
The elementary inequalities are in
\leanmodule{HGAS/Accumulator/RewardGap.lean}.

\subsection{The two cases}

Let \(P\) be the proposition ``\(I\) is solvable'' and choose an ergodic
invariant probability maximizing \(\int\mathcal R\).

If \(P\) holds, ergodicity makes the invariant finite chart have measure
zero or one.  If it had measure one, (11.2) would say that every
conditional drift mean is zero, contradicting (9.6).  It therefore has
measure zero, and (11.3) gives
\[
                    \mathcal M(\mathcal R)\leq-\varepsilon<0.
\tag{11.4}
\]

If \(P\) fails, Section~\ref{sec:no} supplies an invariant graph probability
with reward average at least \(\delta>0\), so
\[
                    \mathcal M(\mathcal R)\geq\delta>0.
\tag{11.5}
\]
Together,
\[
        \boxed{\quad
        \mathcal M(\mathcal R)<0
          \quad\Longleftrightarrow\quad
        I\text{ is solvable}.
        \quad}
\tag{11.6}
\]
This abstract two-case proof is
\leanname{maximalInvariantAverage_neg_iff_of_compactAccumulatorGap} in
\leanmodule{HGAS/Accumulator/CompactSignDichotomy.lean}.

The two instance-dependent facts (9.6) and (10.10) are packaged by
\leanname{gapOfCanonicalThetaOutputs} in
\leanmodule{HGAS/Arithmetic/NormalizedCorrectedEffectiveGapFromCanonicalTheta.lean}.
For the augmented universal family, the resulting unconditional gap is
\leanname{gap} in
\leanmodule{HGAS/Arithmetic/NormalizedCorrectedAugmentedEffectiveFamily.lean}.

\section{The homogeneous product-sphere wrapper}
\label{sec:wrapper}

At this stage the source is a rational polynomial flow tangent to a compact
product of spheres, with reward \(\mathcal R\), and (11.6) is known.  The
wrapper converts this compact angular system into one homogeneous polynomial
field on an ordinary Euclidean space.

\subsection{A rational two-sheeted chart}

Add one coordinate \(z\).  On the Euclidean unit sphere, use the two sheets
\[
                    (k+1)z^2=1,
\tag{12.1}
\]
where \(k\) is the number of source sphere blocks.  Source coordinates are
represented as \(zy\).  If the source field/reward degree bound is the odd
integer \(D\), multiplying chart polynomials by the rational integer
\((k+1)^{(D-1)/2}\) cancels the power of \(z\) on both sheets.  No
\(\sqrt{k+1}\) coefficient is introduced.

Scaled block radii
\[
             s_b=(k+1)\sum_{j\in b}x_j^2
\tag{12.2}
\]
equal \(1\) on the intended full stratum.  Division-free
tangentialization makes the angular code tangent to every block sphere.
A radius-cleanup field drives every other block-radius configuration toward
the full stratum while leaving proper faces invariant.  Proper faces
receive reward \(-1\); on the full stratum the reward is exactly the source
reward \(\mathcal R\).

\subsection{Parity homogenization}

The raw angular and reward polynomials have the correct parity but may
contain several lower degrees.  Multiplying each monomial by a suitable
power of \(\|x\|^2\) makes the angular part homogeneous of
\[
                         d=D+2(k+1)
\tag{12.3}
\]
and the reward homogeneous of degree \(d-1\), without changing either
function on the unit sphere.  The final field is
\[
                           F_{\rm fin}(x)
                  =Y_{\rm hom}(x)+A_{\rm hom}(x)x.
\tag{12.4}
\]
It is rational and homogeneous of odd degree \(d\), and on the sphere its
angular field and radial reward are exactly the wrapped source data.

Invariant-average bookkeeping on the two sheets and proper faces proves
\[
   0\text{ is GAS for }F_{\rm fin}
     \quad\Longleftrightarrow\quad
   \mathcal M(\mathcal R)<0.
\tag{12.5}
\]
The finite construction is
\leanmodule{HGAS/Wrapper/UnitSpherePolynomialWrapper.lean}; the exact GAS
equivalence is
\leanname{CompactSignSourceCertificate.finalField_gas_iff_solvable} in
\leanmodule{HGAS/Compiler/CompactSignSourceWrapper.lean}.

Combining (11.6) and (12.5) gives, for every augmented DQL instance,
\[
       0\text{ is GAS for the emitted homogeneous field}
       \quad\Longleftrightarrow\quad
       I\text{ is solvable}.
\tag{12.6}
\]

\section{Freezing the resources and proving hardness}
\label{sec:fixed}

The remaining question is effectivity.  A semantic equivalence such as
(12.6) is insufficient unless the finite output code is computable from
the universal-halting input and its dimension and degree are independent
of that input.

\subsection{Why the resource bounds are fixed}

In the fixed universal DQL family, input specialization changes
right-hand-side constants but not the quadratic and linear character
tables.  Augmentation preserves this common character data.  Every
compiler degree---the theta detector degree, central clearing exponent,
accumulator degree, wrapper degree, and final homogeneous degree---is a
function of the fixed row/variable counts and common character table.
Therefore it is independent of \(t\).

This is formalized by the
\leanname{SameCharacterData} lemmas in
\leanmodule{HGAS/Arithmetic/NormalizedCorrectedAugmentedUniversalResources.lean}.
The same file fixes:

\begin{itemize}
\item the compact source dimension;
\item the two-block sphere assignment;
\item a proof that the source product sphere is nonempty;
\item the common odd wrapper degree; and
\item the computable finite source field/reward code.
\end{itemize}

\subsection{The effectivity firewall}

The structure
\[
\leanname{FixedUniversalDQLCompactSignSourceFamily}
\]
in
\leanmodule{HGAS/Compiler/FixedUniversalDQLCompactSignSourceFamily.lean}
separates finite code from semantic proof data.  Its parameters and fields
supply:

\begin{enumerate}
\item exact DQL solvability versus unary universal halting;
\item one fixed source block assignment, a nonempty source product sphere,
  and one fixed odd wrapper degree at least three;
\item a compact-sign certificate for every input;
\item one raw output-code function;
\item equality of that raw code with the certificate's field; and
\item computability of the raw output-code function.
\end{enumerate}

The certificate may contain flows, invariant probabilities, real limits,
and classical choices, but the emitted code function is built only from
finite rational syntax.  Its computability theorem composes primitive-
recursive DQL augmentation, controller parameter generation, complex
realification, polynomial compactification, accumulator code, and the
homogeneous wrapper.

For the final family, the theorem
\[
\leanname{hgasQ_output_iff_predicate}
\]
states, for every \(t\),
\[
 \HGAS_{\Q}(n_\ast,d_\ast)
      \bigl(\mathrm{outputCode}(t)\bigr)
 \quad\Longleftrightarrow\quad
 \mathrm{Predicate}(t).
\tag{13.1}
\]
Since \(\mathrm{outputCode}\) is computable, (13.1) is a computable
many-one reduction.  The c.e. upper bound from Theorem~\ref{thm:upper}
then proves
\[
             \HGAS_{\Q}(n_\ast,d_\ast)
             \text{ is }\Sigma^0_1\text{-complete}.
\tag{13.2}
\]
This is precisely
\leanname{fixedPair_sigmaOneComplete}.

\section{Cubicization and cofinal propagation}
\label{sec:cubic}

Write \(d_\ast=2k_\ast+1\), with \(k_\ast\geq1\).  There are two
fixed-parameter computable finite-code transformations.

\subsection{The stabilized Veronese cubic lift}

Map \(x\in\R^{n_\ast}\) to the vector \(v_{k_\ast}(x)\) of all monomials
of degree \(k_\ast\).  If \(F\) has degree \(2k_\ast+1\), differentiating
\(v_{k_\ast}(x(t))\) yields a cubic polynomial expression in the Veronese
coordinates on the Veronese variety.  Away from that variety, add a fixed
rational cubic penalty which drives the toric defining relations to zero.

The formal construction fixes one natural multiplier \(M_{n,k}\) for each
fixed pair \((n,k)\), using \leanname{Classical.choice} on the proved
geometric existence theorem.  For an input field \(F\), the emitted gain
is
\[
       M_{n,k}\bigl(1+\lVert F\rVert_{1,\mathrm{coeff}}\bigr).
\]
The coefficient norm is executable, and
\(\leanname{Primrec.const}\,M_{n,k}\) makes the resulting map primitive
recursive in \(F\).  The formalization therefore proves fixed-\((n,k)\)
computability; it does not formalize real quantifier elimination,
numerically extract \(M_{n,k}\), or give one algorithm computing the
multiplier from varying \((n,k)\).  The manuscript gives a paper-level
uniform strengthening by terminating rational searches with real
quantifier elimination.

The penalty creates a uniformly negative transverse angular block, while
on the toric core the radial cocycle is the source cocycle.  Splitting
invariant probabilities between the core and its complement proves the
exact equivalence
\[
              F\text{ is GAS}
       \quad\Longleftrightarrow\quad
              \operatorname{CubicLift}(F)\text{ is GAS}.
\tag{14.1}
\]
Malformed source codes are sent to a syntactically homogeneous cubic
negative instance, so the map is correct on all raw inputs.

The semantic theorem is
\leanname{sourceGAS_iff_stabilizedCubicGAS} in
\leanmodule{HGAS/Veronese/CubicLiftGASConjugacy.lean}.  The primitive-
recursive sparse compiler and fixed-parameter gain are in
\leanmodule{HGAS/Computability/CubicLiftCodeCompiler.lean} and
\leanmodule{HGAS/Computability/CubicLiftEffectiveGain.lean}.  Its target
dimension is exactly \(N_\ast\) from (1.4).

\subsection{Dimension and odd-degree padding}

From a cubic field, adjoin any fixed number of independent cubic sinks.
To raise degree from \(3\) to \(3+2j\), multiply the old block by a suitable
even radial power and give the new sink coordinates compatible homogeneous
negative dynamics.  Product-GAS and homogeneity calculations give exact
preservation and reflection of GAS.  The finite compiler is
\leanmodule{HGAS/Propagation/DegreePaddingCode.lean}.

For \(N\geq N_\ast\), take \(N-N_\ast\) sink dimensions.  Every odd
\(d\geq3\) has the form
\[
                        d=3+2\frac{d-3}{2}.
\tag{14.2}
\]
Thus the fixed-pair reduction (13.1) propagates to every target pair in
Theorem~\ref{thm:main}.  The exact composition is
\leanname{hgasQ_manyOneReducible_cofinalFixedPair} in
\leanmodule{HGAS/Propagation/CofinalFixedPairReduction.lean}, and
\(\Sigma^0_1\)-membership is supplied again by
Theorem~\ref{thm:upper}.  This completes the proof of
Theorem~\ref{thm:main}.

\section{Relation to the literature}
\label{sec:literature}

\subsection{Stability and decidability}

Classical work on homogeneous systems establishes strong structural
properties and converse Lyapunov theorems; see Samard\v zija
\cite{Samardzija1983} and Rosier~\cite{Rosier1992}.  Ahmadi proved strong
NP-hardness for deciding asymptotic stability of cubic homogeneous vector
fields and explicitly isolated the decidability question
\cite{Ahmadi2012}.  Ahmadi and El Khadir proved algebraic/rational
Lyapunov results for stable homogeneous polynomial fields
\cite{AhmadiElKhadir2020}; Ahmadi and Parrilo studied SOS certificates in
several structured settings~\cite{AhmadiParrilo2017}.

The result here is of a different recursion-theoretic strength.  NP-hardness
does not rule out an algorithm; \(\Sigma^0_1\)-completeness does.  Moreover,
the lower bound already holds for one fixed dimension and degree, and after
the Veronese lift it holds at cubic degree in every sufficiently large fixed
dimension.  The c.e. upper bound identifies the exact arithmetical level and
shows that positive instances can be certified even though no total
decision procedure exists.

There is a broader literature on undecidability and computational
universality for polynomial ODEs, for example Buser and Scarpellini
\cite{BuserScarpellini2007} and Gra\c ca, Campagnolo, and Buescu
\cite{GracaCampagnoloBuescu2008}.  Those results motivate the general
question but do not directly give the fixed homogeneous-GAS statement:
homogeneity, a single equilibrium at the cone tip, forward completeness,
and global attraction impose a substantially different geometry.

\subsection{Arithmetic and special-function inputs}

The arithmetic ancestry is the DPRM theorem
\cite{DavisPutnamRobinson1961,Matiyasevich1970,Matiyasevich1971,
Matiyasevich1993}.  The formal project proves the required primitive-
recursive Diophantine closure internally, then uses the elementary
four-square and multiplication-gate reductions to reach DQL form.

The finite polynomial theta system follows identities standard in the
theory of Jacobi theta functions; useful references are Chandrasekharan
\cite{Chandrasekharan1985}, Mumford~\cite{Mumford1983}, and the NIST DLMF
\cite{DLMF}.  Polynomial dynamical formulations closely related to the
one used here appear in
\cite{BrezhnevLyakhovichSharapov2011,Brezhnev2013}.  The horocycle branch
uses the level-two specialization of the classical periodic-or-
equidistributed dichotomy of Dani--Smillie
\cite{DaniSmillie1984,HubbardMiller2010}.

These citations identify provenance and comparison.  They are not opaque
proof constants in the \Lean{} theorem.  The concrete theta identities,
finite quotient, nonescape estimates, occupation-measure argument,
pointwise equidistribution statement, and every normalization used by the
compiler have formal proofs in the dependency closure.

Kernel checking establishes the stated formal implication; it does not
establish bibliographic priority or replace specialist peer review.  This
is a new manuscript-level result and should be treated as unrefereed until
the reduction and its relationship to the broader literature have received
independent expert review.

\subsection{Scope of the present result}

The formal theorem settles the general and fixed-resource decidability
question for rational homogeneous polynomial fields in the stated
cofinal range.  It leaves at least three natural quantitative problems:

\begin{enumerate}
\item evaluate or drastically reduce the chosen fixed constants;
\item determine the exact low-dimensional boundary, especially
  \(\HGAS_{\Q}(3,3)\); and
\item find a substantially simpler reduction which avoids the modular
  source while preserving fixed cubic resources.
\end{enumerate}

The first two are not consequences of the present proof.  In particular,
the phrase ``cubic undecidability'' throughout this document always means
cubic degree in every sufficiently large fixed dimension.

\section{Lean correspondence and reusable modules}
\label{sec:crosswalk}

The formalization was organized so that the manuscript proof and code have
the same narrow interfaces.  The following table lists the principal
crosswalk.

\begin{center}
\small
\begin{tabular}{>{\raggedright\arraybackslash}p{0.27\textwidth}
                >{\raggedright\arraybackslash}p{0.63\textwidth}}
\toprule
Mathematical step & Principal module or theorem\\
\midrule
Raw rational field language &
\leanmodule{HGAS/Computability/HGASLanguage.lean};
\leanname{HGASQ}\\
\(\Sigma^0_1\) vocabulary and universal source &
\leanmodule{HGAS/Computability/SigmaOne.lean};
\leanname{universalHalts_sigmaOneComplete}\\
C.e. upper bound &
\leanmodule{HGAS/Lyapunov/RationalHomogeneousConverse.lean};
\leanname{hgasQ_re_of_oddDegree}\\
Primitive-recursive Diophantine closure &
\leanmodule{HGAS/Arithmetic/PrimitiveRecursiveDiophantine.lean};
\leanname{binaryPrimrecRel_dioph}\\
Unconditional fixed DQL &
\leanmodule{HGAS/Arithmetic/UnconditionalMRDP.lean} and
\leanmodule{HGAS/Arithmetic/UnconditionalFixedDQL.lean}\\
Nonzero-column repair &
\leanmodule{HGAS/Arithmetic/DQLNonzeroLinearAugmentation.lean};
\leanname{augment_solvable_iff}\\
Positive Fourier selector &
\leanmodule{HGAS/Arithmetic/PositiveConstantTerm.lean} and
\leanmodule{HGAS/Arithmetic/DetectorFourier.lean}\\
Prime-radical nonresonance &
\leanmodule{HGAS/Clocks/RadicalIndependence.lean} and
\leanmodule{HGAS/Arithmetic/NormalizedCorrectedEffectiveSelectedDetectorPhase.lean};
\leanname{physicalDetectorTurnsFrequency_hasNoCharacter}\\
Polynomial Jacobi compiler &
\leanmodule{HGAS/Modular/JacobiPolynomialCode.lean} and
\leanmodule{HGAS/Modular/LinkedJacobiPolynomialCode.lean}\\
Bidirectional recurrent selection &
\leanmodule{HGAS/Arithmetic/NormalizedCorrectedFixedDetectorBidirectionalValidityOrbit.lean}\\
Level-two equidistribution &
\leanmodule{HGAS/Modular/GammaTwoOffPeriodicUnique.lean};
\leanname{gammaTwoDaniSmillieTheorem_unconditional}\\
Canonical root and Bloch continuation &
\leanmodule{HGAS/Arithmetic/NormalizedCorrectedFixedDetectorCanonicalRootSeed.lean} and
\leanmodule{HGAS/Arithmetic/NormalizedCorrectedFixedDetectorPhysicalClassifiedSeed.lean}\\
Repaired three-moment rigidity &
\leanmodule{HGAS/Modular/JacobiThreeMomentTargetRigidity.lean} and
\leanmodule{HGAS/Arithmetic/NormalizedCorrectedEffectiveColumnFrequencyIrrational.lean}\\
Canonical theta output &
\leanmodule{HGAS/Arithmetic/NormalizedCorrectedEffectiveCanonicalThetaOutputs.lean}\\
Yes-instance nonzero drift &
\leanmodule{HGAS/Arithmetic/NormalizedCorrectedEffectiveNonzeroMeanFromTheta.lean}\\
No-instance finite graph &
\leanmodule{HGAS/Arithmetic/NormalizedCorrectedEffectiveNoInstanceFiniteGraph.lean}\\
Compact sign dichotomy &
\leanmodule{HGAS/Accumulator/CompactSignDichotomy.lean} and
\leanmodule{HGAS/Arithmetic/NormalizedCorrectedEffectiveGapFromCanonicalTheta.lean}\\
Homogeneous wrapper &
\leanmodule{HGAS/Compiler/CompactSignSourceWrapper.lean} and
\leanmodule{HGAS/Wrapper/UnitSpherePolynomialWrapper.lean}\\
Fixed effective family &
\leanmodule{HGAS/Compiler/FixedUniversalDQLCompactSignSourceFamily.lean} and
\leanmodule{HGAS/Arithmetic/NormalizedCorrectedAugmentedEffectiveFamily.lean}\\
Cubic/cofinal propagation &
\leanmodule{HGAS/Computability/CubicLiftEffectiveGain.lean} and
\leanmodule{HGAS/Propagation/CofinalFixedPairReduction.lean} and
\leanmodule{HGAS/Propagation/CofinalSigmaOne.lean}\\
\bottomrule
\end{tabular}
\end{center}

Several abstractions are intentionally more reusable than this application:

\begin{itemize}
\item \leanname{DQLNonzeroLinearAugmentation.augment} applies to every
  finite DQL system.
\item \leanname{FixedUniversalDQLCompactSignSourceFamily} accepts any
  fixed-arity DQL family with the required semantic and computable-code
  fields.
\item \leanname{quadraticTrace_eq_target_of_three_moments} is independent
  of theta functions.  Its reusable interface takes quadratic--geometric
  source and target data, nonzero private frequency and twists, an
  irrational ratio, summability, divergent horizons, and three moment
  limits.
\item the invariant-average, ergodic-maximizer, compact-accumulator, and
  product-sphere wrapper theorems do not mention universal halting.
\item the stabilized Veronese GAS reduction accepts every finite
  homogeneous field code in positive dimension and odd degree at least
  three.
\end{itemize}

\section{Trust, audits, and reproducibility}
\label{sec:trust}

\subsection{No hidden mathematical postulate}

The final theorem imports
\leanmodule{HGAS/Arithmetic/NormalizedCorrectedAugmentedEffectiveFamily.lean}
through \leanmodule{HGAS.lean}.  The source tree contains no declaration
matching a strict scan for project \leanname{axiom}, \leanname{unsafe},
\leanname{sorry}, or \leanname{admit}.  In particular:

\begin{itemize}
\item MRDP closure is a theorem, not an axiom;
\item the level-two Dani--Smillie statement is a theorem, not an interface
  hypothesis;
\item canonical theta outputs are proved for the augmented family;
\item the no-instance graph is constructed, not assumed; and
\item the final family has no remaining gap or theta parameter.
\end{itemize}

Running \texttt{\#print axioms} on
\leanname{fixedPair_sigmaOneComplete},
\leanname{cofinal_sigmaOneComplete},
\leanname{cubic_sigmaOneComplete},
\leanname{fixedPair_not_computable}, and
\leanname{cubic_not_computable} (as well as the two named
non-c.e.-complement consequences) returns exactly
\[
       [\,\leanname{propext},\
          \leanname{Classical.choice},\
          \leanname{Quot.sound}\,].
\tag{16.1}
\]
These are ordinary logical foundations of \Lean{}/\Mathlib{}; there is no
project-specific axiom and no \leanname{sorryAx}.

The audited environment is Lean \(4.30.0\) with Mathlib tag \(v4.30.0\),
commit
\leanname{c5ea00351c28e24afc9f0f84379aa41082b1188f}.
For the prover and library, see de Moura and Ullrich~\cite{Lean4} and the
Mathlib community~\cite{Mathlib}.

\subsection{Mechanical checks}

From the \leanmodule{lean/} directory, the decisive commands are:

\begin{verbatim}
lake build HGAS

lake env lean -DwarningAsError=true \
  HGAS/Arithmetic/NormalizedCorrectedAugmentedEffectiveFamily.lean

rg -n \
  -e '^[[:space:]]*((private|protected)[[:space:]]+)?axiom[[:space:]]' \
  -e '^[[:space:]]*((private|protected)[[:space:]]+)?unsafe[[:space:]]' \
  -e '^[[:space:]]*(sorry|admit)\b' \
  -e ':=[[:space:]]*(by[[:space:]]+)?(sorry|admit)\b' \
  HGAS HGAS.lean
\end{verbatim}

The public/default \leanname{HGAS} target build completes successfully; the
endpoint file also recompiles successfully with warnings treated as errors,
and the strict token scan is empty.  Its transitive import closure contains
783 of the 824 modules under \leanmodule{HGAS/}.  All 41 modules outside
that closure were also elaborated directly with
\texttt{-DwarningAsError=true}; all succeeded.  Their exact list and
legacy/audit status are recorded in \leanmodule{LEGACY_MODULES.md}.
Independent audit reports are stored under \leanmodule{../audits/} from
this working directory.  They separately check the MRDP chain, theta
identities, clock normalization, quotient/horocycle arguments, selected
phase factorization, compact accumulator, wrapper, cubic lift, augmented
fixed resources, and the repaired endpoint.

\subsection{Why classical choice does not spoil the reduction}

The fixed MRDP polynomial and some analytic witnesses, including the
geometry-only cubic gain multiplier for each fixed \((n,k)\), are obtained
by classical choice.  This is compatible with a fixed-parameter computable
many-one reduction because each chosen finite object is thereafter a
constant.  A program may contain any fixed finite lookup table.  This does
not assert a uniform extractor for those constants as the parameters vary.
More importantly, the final
family structure does not derive its output code by inspecting semantic
proof objects: it stores a separate finite source-code function and a proof
that this function is computable.  The equality with the semantic
certificate is proved afterward.  This is the formal effectivity firewall
described in Section~\ref{sec:fixed}.

\subsection{Final logical compression}

After all mathematics has been packaged, the endpoint proof is short:
\[
\begin{array}{rcl}
\mathrm{Predicate}(t)
 &\Longleftrightarrow&
 \widetilde I_t\text{ solvable}\\
 &\Longleftrightarrow&
 \mathcal M(\mathcal R_t)<0\\
 &\Longleftrightarrow&
 \mathrm{outputCode}(t)\in
       \HGAS_{\Q}(n_\ast,d_\ast).
\end{array}
\tag{16.2}
\]
The first equivalence is arithmetic and augmentation; the second is the
two-field semantic gap; the third is the homogeneous wrapper.  Computability
of \(\mathrm{outputCode}\) makes (16.2) a many-one reduction.  Combining it
with the c.e. upper bound proves fixed-pair \(\Sigma^0_1\)-completeness;
the effective Veronese and padding maps give the cofinal theorem.

\section{Conclusion}

For rational homogeneous polynomial vector fields, the general GAS
decision problem is not merely difficult: it is undecidable already at one
fixed dimension and one fixed odd degree.  Positive instances remain
computably enumerable through exact rational Lyapunov certificates, so the
fixed language is precisely \(\Sigma^0_1\)-complete.  A stabilized
Veronese lift lowers the degree to three at the cost of a fixed increase in
dimension, and padding propagates the result to every sufficiently large
fixed dimension and every odd degree at least three.

The proof's delicate point is also its most instructive correction.  Three
moments do not determine a quadratic theta trace when the linear/private
frequency ratio is zero.  A tiny arithmetic preprocessing step,
\(x_j-y_j=0\), both preserves solvability and makes the ratio irrational
through prime-radical character independence.  The repaired theorem is
therefore stronger logically and cleaner modularly: every final universal
instance passes through one uniform, reusable irrational-ratio selector.

The low-dimensional cubic case \((n,d)=(3,3)\) remains open in this
development.  Settling it would be a separate result rather than a missing
step in the theorem proved here.

\begin{thebibliography}{99}
\small
\setlength{\itemsep}{1pt}

\bibitem{Ahmadi2012}
A.~A. Ahmadi,
\newblock On the difficulty of deciding asymptotic stability of cubic
homogeneous vector fields,
\newblock in \emph{Proceedings of the 2012 American Control Conference},
2012, pp.~3334--3339,
\newblock \href{https://doi.org/10.1109/ACC.2012.6315687}
{doi:10.1109/ACC.2012.6315687};
\href{https://arxiv.org/abs/1112.0741}{arXiv:1112.0741}.

\bibitem{AhmadiElKhadir2020}
A.~A. Ahmadi and B. El Khadir,
\newblock On algebraic proofs of stability for homogeneous vector fields,
\newblock \emph{IEEE Transactions on Automatic Control} 65 (2020), no.~1,
325--332,
\newblock \href{https://doi.org/10.1109/TAC.2019.2914968}
{doi:10.1109/TAC.2019.2914968};
\href{https://arxiv.org/abs/1803.01877}{arXiv:1803.01877}.

\bibitem{AhmadiParrilo2017}
A.~A. Ahmadi and P.~A. Parrilo,
\newblock Sum of squares certificates for stability of planar,
homogeneous, and switched systems,
\newblock \emph{IEEE Transactions on Automatic Control} 62 (2017), no.~10,
5269--5274,
\newblock \href{https://doi.org/10.1109/TAC.2016.2647253}
{doi:10.1109/TAC.2016.2647253}.

\bibitem{Brezhnev2013}
Y.~V. Brezhnev,
\newblock Non-canonical extension of theta-functions and modular
integrability of theta-constants,
\newblock \emph{Proceedings of the Royal Society of Edinburgh Section A:
Mathematics} 143 (2013), no.~4, 689--738,
\newblock \href{https://doi.org/10.1017/S0308210512001023}
{doi:10.1017/S0308210512001023};
\href{https://arxiv.org/abs/1011.1643}{arXiv:1011.1643}.

\bibitem{BrezhnevLyakhovichSharapov2011}
Y.~V. Brezhnev, S.~L. Lyakhovich, and A.~A. Sharapov,
\newblock Dynamical systems defining Jacobi's theta-constants,
\newblock \emph{Journal of Mathematical Physics} 52 (2011), no.~11,
112704,
\newblock \href{https://doi.org/10.1063/1.3662961}
{doi:10.1063/1.3662961};
\href{https://arxiv.org/abs/1012.1429}{arXiv:1012.1429}.

\bibitem{BuserScarpellini2007}
P. Buser and B. Scarpellini,
\newblock Undecidable propositions by ODE's,
\newblock \emph{Annales Academiae Scientiarum Fennicae Mathematica}
32 (2007), 317--340,
\newblock
\href{https://www.acadsci.fi/mathematica/Vol32/vol32pp317-340.pdf}
{journal PDF}.

\bibitem{Chandrasekharan1985}
K. Chandrasekharan,
\newblock \emph{Elliptic Functions},
\newblock Grundlehren der mathematischen Wissenschaften 281,
Springer, Berlin, 1985,
\newblock \href{https://doi.org/10.1007/978-3-642-52244-4}
{doi:10.1007/978-3-642-52244-4}.

\bibitem{Dani1978}
S.~G. Dani,
\newblock Invariant measures of horospherical flows on noncompact
homogeneous spaces,
\newblock \emph{Inventiones Mathematicae} 47 (1978), no.~2, 101--138,
\newblock \href{https://doi.org/10.1007/BF01578067}
{doi:10.1007/BF01578067}.

\bibitem{DaniSmillie1984}
S.~G. Dani and J. Smillie,
\newblock Uniform distribution of horocycle orbits for Fuchsian groups,
\newblock \emph{Duke Mathematical Journal} 51 (1984), no.~1, 185--194,
\newblock \href{https://doi.org/10.1215/S0012-7094-84-05110-X}
{doi:10.1215/S0012-7094-84-05110-X}.

\bibitem{DavisPutnamRobinson1961}
M. Davis, H. Putnam, and J. Robinson,
\newblock The decision problem for exponential Diophantine equations,
\newblock \emph{Annals of Mathematics}, Second Series, 74 (1961), no.~3,
425--436,
\newblock \href{https://doi.org/10.2307/1970289}
{doi:10.2307/1970289}.

\bibitem{DLMF}
F.~W.~J. Olver, A.~B. Olde Daalhuis, D.~W. Lozier, B.~I. Schneider,
R.~F. Boisvert, C.~W. Clark, B.~R. Miller, B.~V. Saunders, H.~S. Cohl,
and M.~A. McClain (eds.),
\newblock \emph{NIST Digital Library of Mathematical Functions},
\newblock version 1.2.7, released 15 June 2026,
\url{https://dlmf.nist.gov/}; especially Chapters 20 and 23.

\bibitem{GracaCampagnoloBuescu2008}
D.~S. Gra\c ca, M.~L. Campagnolo, and J. Buescu,
\newblock Computability with polynomial differential equations,
\newblock \emph{Advances in Applied Mathematics} 40 (2008), no.~3,
330--349,
\newblock \href{https://doi.org/10.1016/j.aam.2007.02.003}
{doi:10.1016/j.aam.2007.02.003}.

\bibitem{HubbardMiller2010}
J.~H. Hubbard and R.~L. Miller,
\newblock Equidistribution of horocyclic flows on complete hyperbolic
surfaces of finite area,
\newblock in \emph{In the Tradition of Ahlfors--Bers, V},
Contemporary Mathematics 510, American Mathematical Society, 2010,
pp.~129--150,
\newblock \href{https://doi.org/10.1090/conm/510/10022}
{doi:10.1090/conm/510/10022};
\href{https://arxiv.org/abs/0712.1300}{arXiv:0712.1300}.

\bibitem{Lean4}
L. de Moura and S. Ullrich,
\newblock The Lean 4 theorem prover and programming language,
\newblock in \emph{Automated Deduction---CADE 28}, Lecture Notes in
Computer Science 12699, Springer, 2021, pp.~625--635,
\newblock \href{https://doi.org/10.1007/978-3-030-79876-5_37}
{doi:10.1007/978-3-030-79876-5\_37}.

\bibitem{Lagrange1770}
J.-L. Lagrange,
\newblock D\'emonstration d'un th\'eor\`eme d'arithm\'etique,
\newblock \emph{Nouveaux M\'emoires de l'Acad\'emie Royale des Sciences et
Belles-Lettres de Berlin}, volume for 1770 (printed 1772), 123--133.

\bibitem{Matiyasevich1970}
Y.~V. Matiyasevich,
\newblock Enumerable sets are Diophantine,
\newblock \emph{Soviet Mathematics Doklady} 11 (1970), no.~2, 354--357;
English translation of \emph{Doklady Akademii Nauk SSSR} 191 (1970),
no.~2, 279--282.

\bibitem{Matiyasevich1971}
Y.~V. Matiyasevich,
\newblock Diophantine representation of enumerable predicates,
\newblock \emph{Mathematics of the USSR-Izvestiya} 5 (1971), no.~1,
1--28,
\newblock \href{https://doi.org/10.1070/IM1971v005n01ABEH001004}
{doi:10.1070/IM1971v005n01ABEH001004}.

\bibitem{Matiyasevich1993}
Y.~V. Matiyasevich,
\newblock \emph{Hilbert's Tenth Problem},
\newblock Foundations of Computing, MIT Press, Cambridge, MA, 1993.

\bibitem{Mathlib}
The Mathlib Community,
\newblock The Lean mathematical library,
\newblock in \emph{Proceedings of the 9th ACM SIGPLAN International
Conference on Certified Programs and Proofs (CPP 2020)}, 367--381,
\newblock \href{https://doi.org/10.1145/3372885.3373824}
{doi:10.1145/3372885.3373824};
\href{https://arxiv.org/abs/1910.09336}{arXiv:1910.09336}.

\bibitem{Mumford1983}
D. Mumford,
\newblock \emph{Tata Lectures on Theta I},
\newblock Progress in Mathematics 28, Birkh\"auser, Boston, 1983,
\newblock \href{https://doi.org/10.1007/978-1-4899-2843-6}
{doi:10.1007/978-1-4899-2843-6}.

\bibitem{Rosier1992}
L. Rosier,
\newblock Homogeneous Lyapunov function for homogeneous continuous vector
field,
\newblock \emph{Systems \& Control Letters} 19 (1992), no.~6, 467--473,
\newblock \href{https://doi.org/10.1016/0167-6911(92)90078-7}
{doi:10.1016/0167-6911(92)90078-7}.

\bibitem{Samardzija1983}
N. Samard{\v z}ija,
\newblock Stability properties of autonomous homogeneous polynomial
differential systems,
\newblock \emph{Journal of Differential Equations} 48 (1983), no.~1,
60--70,
\newblock \href{https://doi.org/10.1016/0022-0396(83)90059-1}
{doi:10.1016/0022-0396(83)90059-1}.

\end{thebibliography}

\end{document}
